\section*{Chapter \ref{ch:b2:ellingham} --- Ellingham Diagrams and Metallurgy}

\begin{solution}{exo:b2:ellingham:1}
\ce{4Cu + O2 -> 2Cu2O}; \ce{4/3Al + O2 -> 2/3Al2O3}; \ce{2C + O2 -> 2CO};
\ce{C + O2 -> CO2}.
\end{solution}

\begin{solution}{exo:b2:ellingham:2}
$\Delta_r H^\circ = 2(-350.46) = \qty{-700.92}{kJ}$; $\Delta_r S^\circ = 2(43.65)
- 2(41.63) - 205.15 = \qty{-201.1}{J/K}$. So $\Delta_r G^\circ = -700.9 +
0.2011\,T$ (\unit{kJ}, $T$ in \unit{K}), valid up to the melting point of zinc.
\end{solution}

\begin{solution}{exo:b2:ellingham:3}
At \qty{1500}{K} the line of \ce{Cr2O3} is near \qty{-495}{kJ}. Below it, so
able to reduce chromium oxide: silicon, titanium, aluminium, magnesium,
calcium. Above it, unable: copper, iron, zinc. Carbon, at \qty{-488}{kJ}, lies
just above: it reduces chromium oxide only at a slightly higher temperature,
and then gives a chromium loaded with carbide --- one reason why chromium for
alloys free of carbon is made by aluminothermy.
\end{solution}

\begin{solution}{exo:b2:ellingham:4}
The slope is $-\Delta_r S^\circ$, governed by the change in the amount of gas.
\ce{C + O2 -> CO2}: one mole of gas gives one, $\Delta_r S^\circ \approx 0$,
horizontal. \ce{2C + O2 -> 2CO}: one gives two, $\Delta_r S^\circ > 0$, the
line falls. \ce{2CO + O2 -> 2CO2}: three give two, $\Delta_r S^\circ < 0$, the
line rises, like those of the metals.
\end{solution}

\begin{solution}{exo:b2:ellingham:5}
$\Delta_r G^\circ(\qty{600}{K}) = -181.6 + 600 \times 0.2165 = \qty{-51.7}{kJ}$,
$p_{\ce{O2}} = p^\circ\exp(-51\,710/(8.314 \times 600)) = \qty{3.2e-5}{bar}$. Air
holds \qty{0.21}{bar}, much more: mercury is oxidised at \qty{600}{K} (slowly
--- this is how the oxide was first prepared), and the oxide decomposes only
above about \qty{730}{K}.
\end{solution}

\begin{solution}{exo:b2:ellingham:6}
$\Delta_r G^\circ = -1582.3 - (-743.5) = \qty{-838.8}{kJ}$. The reagents are
two solids in contact only at the surface of their grains, and the reaction
has a high activation energy: it must be ignited locally (a magnesium fuse);
then the heat it releases keeps it going.
\end{solution}

\begin{solution}{exo:b2:ellingham:7}
The line of \ce{2C + O2 -> 2CO} crosses that of \ce{2Fe + O2 -> 2FeO} near
\qty{1040}{K}; above, \ce{FeO + C -> Fe + CO} has a negative \omterm{def:b2:reaction-free-energy:reaction-gibbs}{standard reaction
Gibbs energy}.
\end{solution}

\begin{solution}{exo:b2:ellingham:8}
$\Delta_r G^\circ = 2(-209.1) - (-396.0) = \qty{-22.2}{kJ}$, $K^\circ =
\exp(22\,200/(8.314 \times 1100)) = 11.3$. Then $x^2/(1 - x) = 11.3$, $x =
0.92$. On cooling to \qty{700}{K} the equilibrium moves back towards \ce{CO2}
($x = 0.016$): carbon monoxide disproportionates, \ce{2CO -> C + CO2}, and
soot deposits --- slowly, unless a surface such as iron catalyses it.
\end{solution}

\begin{solution}{exo:b2:ellingham:9}
The water line (three moles of gas give two) rises less steeply than the
\ce{CO}/\ce{CO2} line, also three for two but with a larger entropy loss. The
two cross near \qty{1100}{K}: below, the \ce{CO}/\ce{CO2} line is lower and
carbon monoxide is the stronger reductant; above, the water line is lower and
hydrogen wins. It is the same statement as the shift of the water-gas
equilibrium \ce{CO + H2O <=> CO2 + H2} towards the left at high temperature.
\end{solution}

\begin{solution}{exo:b2:ellingham:10}
Per mole of \ce{O2}, that is for \ce{2MgO + Si -> SiO2 + 2Mg(g)}: $\Delta_r
G^\circ = -643.7 - (-845.5) = \qty{+201.8}{kJ}$. With silicon and silica of
activity 1, $\Delta_r G = \Delta_r G^\circ + RT\ln(p_{\ce{Mg}}/p^\circ)^2$, negative
when $p_{\ce{Mg}} < p^\circ\exp(-201\,800/(2 \times 8.314 \times 1500)) =
\qty{3e-4}{bar}$. A vacuum furnace pumps the magnesium vapour away as fast
as it forms and condenses it on a cold surface, so its pressure never reaches
that value; lime added to the charge binds the silica and lowers its activity,
which helps further.
\end{solution}

\begin{solution}{exo:b2:ellingham:11}
Four species, one reaction, three phases (two solids and the gas): $v = 4 - 1 +
2 - 3 = 2$. At given temperature and pressure the ratio $p_{\ce{CO}}/p_{\ce{CO2}}$
is fixed: the gas cannot give up all its carbon monoxide to the iron oxide,
and the top gas still contains much of it --- which is burnt to preheat the
air.
\end{solution}

\begin{solution}{exo:b2:ellingham:12}
$\log K = 2(0.34 + 0.41)/0.059 = 25.4$, $K = \num{3e25}$. With
$[\ce{Fe^2+}] = \qty{1.0}{mol/L}$, $[\ce{Cu^2+}] = 1/K \approx
\qty{4e-26}{mol/L}$: the copper is removed completely.
\end{solution}

\begin{solution}{pb:b2:ellingham:1}
\textbf{1.} $\Delta_r H^\circ = \qty{-700.92}{kJ}$, $\Delta_r S^\circ =
\qty{-201.1}{J/K}$.
\textbf{2.} $\Delta_r G^\circ = -700.9 + 0.2011\,T$ (\unit{kJ}), below
\qty{692.7}{K}.
\textbf{3.} $\Delta_{\text{fus}}S = 7320/692.7 = \qty{10.6}{J/(K.mol)}$;
$\Delta_{\text{vap}}S = 115\,300/1180.2 = \qty{97.7}{J/(K.mol)}$.
\textbf{4.} Liquid zinc: $\Delta_r H^\circ = -700.92 - 2(7.32) = \qty{-715.6}{kJ}$,
$\Delta_r S^\circ = -201.1 - 2(10.6) = \qty{-222.2}{J/K}$: $\Delta_r G^\circ =
-715.6 + 0.2222\,T$.
\textbf{5.} Gaseous zinc: $\Delta_r H^\circ = -715.6 - 2(115.3) = \qty{-946.2}{kJ}$,
$\Delta_r S^\circ = -222.2 - 2(97.7) = \qty{-417.7}{J/K}$: $\Delta_r G^\circ =
-946.2 + 0.4177\,T$.
\textbf{6.} Slopes $0.201$, $0.222$ and $0.418$ \unit{kJ/K}. Above the boiling
point the reaction consumes three moles of gas instead of one: the entropy
loss, hence the slope, doubles.
\textbf{7.} At \qty{692.7}{K}: $-700.9 + 139.3 = -715.6 + 153.9 = \qty{-561.6}{kJ}$;
at \qty{1180.2}{K}: $-715.6 + 262.3 = -946.2 + 492.9 = \qty{-453.3}{kJ}$. The
lines join, as they must.
\textbf{8.} Slope $(-453.0 + 418.2)/200 = \qty{-0.174}{kJ/K}$, $a = -418.2 +
0.174 \times 1100 = \qty{-226.8}{kJ}$: $\Delta_r G^\circ = -226.8 - 0.174\,T$.
\textbf{9.} $\Delta_r S^\circ = \qty{+174}{J/K}$: one mole of gas gives two.
\textbf{10.} \ce{2ZnO + 2C -> 2Zn(g) + 2CO}, $\Delta_r G^\circ = (-226.8 - 0.174\,T)
- (-946.2 + 0.4177\,T) = 719.4 - 0.592\,T$ (\unit{kJ}).
\textbf{11.} Zero at $T = 719.4/0.592 = \qty{1215}{K}$.
\textbf{12.} Above its boiling point zinc forms as a vapour: the reaction makes
gas from solids, its entropy is large and positive, and the vapour leaves the
charge, which separates the metal from the ore and the coke by distillation.
\textbf{13.} Four species, one reaction, one particular condition ($p_{\ce{Zn}} =
p_{\ce{CO}}$), three phases: $v = 4 - 1 - 1 + 2 - 3 = 1$. Fixing the pressure
fixes the temperature of equilibrium.
\textbf{14.} One mole of zinc vapour per mole of carbon monoxide:
$p_{\ce{Zn}} = p_{\ce{CO}} = \qty{0.5}{bar}$.
\textbf{15.} Per mole of \ce{ZnO}, $\Delta_r G^\circ = \frac12(719.4 - 0.592 \times
1300) = \qty{-25.1}{kJ/mol}$, $K^\circ = \exp(25\,100/(8.314 \times 1300)) = 10$;
$Q = 0.25 < K^\circ$: the reduction runs.
\textbf{16.} Just above the crossing the reduction is already favoured (with
\qty{0.5}{bar} of each gas, from about \qty{1170}{K}); hotter costs fuel,
wears the clay retorts and does not make more zinc, since the reaction goes to
completion anyway.
\textbf{17.} \ce{Zn(g) + CO2 -> ZnO + CO}.
\textbf{18.} The line of zinc oxide lies below the \ce{CO}/\ce{CO2} line
($-445$ against \qty{-357}{kJ} at \qty{1200}{K}): zinc vapour reduces carbon
dioxide and turns back into oxide. In the retort the hot carbon destroys any
\ce{CO2} (Boudouard); in the cooler condenser nothing does, and \ce{CO2} --- from
air leaking in, or from \ce{2CO -> C + CO2} on cooling --- spoils the zinc into a
grey oxidised powder.
\textbf{19.} Slow cooling leaves the vapour for a long time in the range where
it is reoxidised, and gives a fine dust instead of liquid metal; quick
condensation to the liquid limits both.
\textbf{20.} $10^6/65.38 = \qty{1.53e4}{mol}$ of zinc, as many of carbon:
$1.53 \times 10^4 \times 12.011 = \qty{184}{kg}$ (much more in practice, to heat the
retorts).
\textbf{21.} $\Delta_r H^\circ = \frac12(-226.8 + 946.2) = \qty{+360}{kJ}$ per
mole of zinc; per tonne, $1.53 \times 10^4 \times 360 = \qty{5.5e6}{kJ}$, that is
\qty{5.5}{GJ} --- supplied by burning more coal around the retorts.
\textbf{22.} Carbon reduces zinc oxide under standard conditions above
$\boldsymbol{T \approx \qty{1.22e3}{K}}$.
\end{solution}
